\documentclass[11pt,a4paper]{article}
\usepackage[margin=1.8cm]{geometry}
\usepackage{amsmath}
\usepackage{booktabs}
\usepackage{graphicx}
\usepackage{hyperref}
\usepackage{xcolor}
\usepackage{listings}

\setlength{\parskip}{0.35em}
\setlength{\parindent}{0pt}
\setlength{\tabcolsep}{4pt}
\renewcommand{\arraystretch}{0.95}

\definecolor{codebg}{rgb}{0.95,0.95,0.95}
\lstset{
    backgroundcolor=\color{codebg},
    basicstyle=\ttfamily\footnotesize,
    breaklines=true,
    frame=single,
    tabsize=4
}

\title{\vspace{-1.2cm}\textbf{Furuta Pendulum Sim-to-Real Report}\\
\large Current MuJoCo Model, Hardware-Mass Training, and Domain Randomization}
\author{}
\date{June 9, 2026}

\begin{document}
\maketitle
\vspace{-1.0cm}

\section{Summary}

This project trains a PPO controller for a Furuta pendulum in MuJoCo before
transfer to real hardware.  The model uses SolidWorks CAD meshes converted to
MuJoCo, with the arm angle limited to $\pm135^\circ$ so the mechanism does not
spin continuously, wrap cables, or block the future CyberRunner camera view.

The most recent update replaces the older CAD-estimated masses with measured
hardware masses.  A new clean policy and a new domain-randomized (DR) policy
were trained from these masses.  The current deployment candidate is:

\begin{center}
\texttt{trained\_policies/dr\_hardware\_masses\_20260609/best\_model.zip}
\end{center}

\begin{center}
\begin{tabular}{llll}
\toprule
\textbf{Policy} & \textbf{Run} & \textbf{Best mean} & \textbf{Max reward} \\
\midrule
Clean hardware-mass policy & \texttt{20260609\_133511} & 4095.68 & 4200 \\
Clean hardware-mass policy & \texttt{20260609\_133618} & 4094.71 & 4200 \\
DR hardware-mass policy & \texttt{20260609\_135537} & 4051.84 & 4200 \\
\bottomrule
\end{tabular}
\end{center}

The latest completed experiment adds motor deadband and elbow velocity
filtering.  It warm-started from the hardware-mass DR policy and better matches
real motor and AS5600 encoder behavior.

\section{System Model}

The Furuta pendulum has a motorized rotary arm and a passive pendulum attached
at the arm tip.  The motor does not directly actuate the pendulum; instead, it
moves the arm to accelerate the pendulum pivot.  During swing-up, the controller
injects energy into the pendulum.  During balance, it makes corrective arm
motions to keep the pendulum upright.

\begin{center}
\begin{tabular}{ll}
\toprule
\textbf{Item} & \textbf{Current value} \\
\midrule
Motor & 25GA20E260, 7.4 V DC, 1:20 gearbox \\
Pendulum sensor & AS5600, 12-bit magnetic encoder \\
CAD files & \texttt{fixed\_motor.stl}, \texttt{arm\_2.STL}, \texttt{rod\_2.STL} \\
Mesh scale & \texttt{0.001 0.001 0.001} (mm to m) \\
Mesh refpos & \texttt{602.303 877.409 1307.925} mm \\
Coordinate remap & SolidWorks X-up to MuJoCo Z-up via \texttt{xyaxes="0 0 1 1 0 0"} \\
Control period & 10 ms \\
\bottomrule
\end{tabular}
\end{center}

The angle convention is:
\[
\theta=0 \text{ upright}, \qquad \theta=\pm\pi \text{ hanging down},
\]
where $\theta$ is the pendulum angle.  The shoulder/arm angle is $\phi$.

\section{MuJoCo Parameters}

\begin{center}
\begin{tabular}{lll}
\toprule
\textbf{Joint} & \textbf{Definition} & \textbf{Notes} \\
\midrule
Shoulder & \texttt{pos="0 0 0.02575" axis="0 0 1"} & actuated, $\pm135^\circ$, damping 0.01 \\
Elbow & \texttt{pos="0 -0.049505 0.056747" axis="0 1 0" ref="0"} & passive, $0$ rad upright \\
\bottomrule
\end{tabular}
\end{center}

\begin{center}
\begin{tabular}{lrr}
\toprule
\textbf{Body} & \textbf{Measured mass (g)} & \textbf{XML value (kg)} \\
\midrule
fixed\_motor & 92.00 & 0.09200 \\
arm & 43.00 & 0.04300 \\
pendulum\_rod & 15.00 & 0.01500 \\
\bottomrule
\end{tabular}
\end{center}

The actuator is modeled as a normalized shoulder torque command:
\[
u\in[-1,1] \quad \rightarrow \quad \tau\in[-0.132,0.132]\ \mathrm{Nm}.
\]

\section{PPO Controller}

The policy observation is:
\[
\mathbf{o}=[\cos\theta,\ \sin\theta,\ \dot{\theta},\ \phi,\ \dot{\phi}].
\]
Using $\cos\theta$ and $\sin\theta$ avoids the discontinuity at $\pm\pi$.
The action is the normalized motor command $u$.

PPO was selected because it is stable, simple to implement with
Stable-Baselines3, and works well for continuous-control MuJoCo tasks.  PPO is
an on-policy actor-critic method: it collects rollout data using the current
policy, estimates advantages, then updates the policy using a clipped objective
so each update does not move too far from the behavior that generated the data.

The reward used for the current hardware-mass training is:
\[
r=\cos(\theta)-0.1u^2
-0.5\max(0,\ |\phi|-2.094)^2
+0.2\mathbf{1}_{|\theta|<10^\circ}
+0.2\mathbf{1}_{|\theta|<5^\circ}
-0.005\dot{\theta}^2\mathbf{1}_{|\theta|<10^\circ}.
\]

The terms reward upright balance, discourage excessive torque, discourage the
arm from approaching its mechanical limit, reward tighter upright capture, and
penalize high pendulum velocity once near upright.  A 30 s episode has 3000
control steps, so the maximum score is approximately:
\[
3000\times1.4=4200.
\]

\section{Domain Randomization}

Domain randomization is used so the policy does not depend on a perfect
simulation.  The policy is first trained in clean simulation, then fine-tuned
gradually with randomized parameters.  This curriculum is important because
training directly under all imperfections can make PPO forget the clean swing-up
and balance behavior before it learns robustness.

\begin{center}
\begin{tabular}{lll}
\toprule
\textbf{Parameter} & \textbf{Current range} & \textbf{Purpose} \\
\midrule
Mass/inertia & $\pm5\%$ & measurement uncertainty \\
Shoulder damping & $\pm5\%$ & gearbox/friction variation \\
Elbow damping & $[0,\ 0.00002]$ Nm s/rad & tiny pivot friction \\
Motor gear/torque & $\pm5\%$ & voltage/motor variation \\
Action delay & 0--1 step & timing margin \\
Shoulder position noise & $\sigma=0.017$ rad & backlash / encoder error \\
Shoulder velocity noise & $\sigma=0.050$ rad/s & velocity estimate uncertainty \\
Elbow position noise & $\sigma=0.001$ rad & AS5600 angle noise \\
Shoulder zero offset & $\pm1^\circ$ & calibration tolerance \\
Elbow zero offset & $\pm0.5^\circ$ & upright zero tolerance \\
\bottomrule
\end{tabular}
\end{center}

\subsection{Deadband and Filter Experiment}

The latest training run adds two additional hardware effects:
\begin{itemize}
    \item \textbf{Motor deadband:} commands below a sampled threshold produce no
          torque.  The current gentle range is 0--5\% of full command.
    \item \textbf{Elbow velocity filtering:} the AS5600 measures angle, so
          velocity must be estimated by differentiation and low-pass filtering.
          The current filter range is $\alpha\in[0.7,0.9]$.
\end{itemize}

The run is:
\[
\texttt{runs/dr/20260609\_deadband5\_filter70\_90}.
\]
It warm-started from \texttt{runs/dr/20260609\_hardware\_masses}.  Early
evaluation dropped to roughly 1100--1900 / 4200 because the new motor deadband
and velocity lag made balance harder.  After additional PPO updates, the best
checkpoint reached 4032.66 / 4200 with standard deviation 31.09 and full
3000-step episodes.  A recorded deterministic GIF from this checkpoint scored
4049.2 / 4200.

\section{Current Results}

\begin{center}
\begin{tabular}{llll}
\toprule
\textbf{Stage} & \textbf{Run} & \textbf{Result} & \textbf{Comment} \\
\midrule
Clean, old masses & \texttt{20260604\_163108} & 2886.57 / 3000 & historical baseline \\
DR, old masses & \texttt{20260604\_204533} & 2554.50 / 3000 & historical robustness policy \\
Long hold + kicks, old masses & \texttt{20260605\_elbow\_kick\_0p5deg\_x3} & 8233.1 / 8400 & historical disturbance policy \\
Clean, hardware masses & \texttt{20260609\_133511} & 4095.68 / 4200 & new clean baseline \\
Clean, hardware masses & \texttt{20260609\_133618} & 4094.71 / 4200 & repeat clean baseline \\
DR, hardware masses & \texttt{20260609\_135537} & 4051.84 / 4200 & current deployment candidate \\
DR + deadband/filter & \texttt{20260609\_deadband5\_filter70\_90} & 4032.66 / 4200 & best hardware-effect policy \\
\bottomrule
\end{tabular}
\end{center}

\section{Next Steps}

\begin{enumerate}
    \item Test the hardware-mass DR policy on the real pendulum after AS5600
          upright zero calibration.
    \item Test the deadband/filter DR policy if the hardware-mass policy is too
          idealized for real motor or sensor behavior.
    \item If the deadband/filter policy is still too sensitive on hardware,
          reduce the added hardware effects further, for example 0--2.5\%
          motor deadband or $\alpha\in[0.8,0.95]$.
    \item Identify real system parameters and update the MuJoCo model / DR
          ranges based on measurements.
    \item Model the tilting board, place the pendulum on it, and train the next
          simulation phase with board tilt disturbances.
\end{enumerate}

\begin{thebibliography}{9}
\bibitem{schulman2017proximal}
J.~Schulman, F.~Wolski, P.~Dhariwal, A.~Radford, and O.~Klimov,
``Proximal Policy Optimization Algorithms,'' \textit{arXiv:1707.06347}, 2017.
\end{thebibliography}

\end{document}
